\section{Matrices over a commutative semiring}
\label{sec:matrix}

Suppose $S = (S, +, 0, \cdot, 1)$ a commutative semiring. Write $\sum_i a_i$ for $a_1 + \cdots + a_n$. Two
derived properties of $\sum$ will be useful: \emph{distributivity} of $\cdot$ over $\sum$, $a \cdot \sum_i b_i
= \sum_i a \cdot b_i$ (and symmetrically on the right), which follows from the distributivity of $\cdot$ over
$+$, and the \emph{interchange law} $\sum_i \sum_j a_{ij} = \sum_j \sum_i a_{ij}$, which follows from
associativity and commutativity of $+$.

\subsection{Vectors, matrices, and standard basis}
\label{sec:matrix:vectors}

An \emph{$n$-length vector over $X$} is an element $v \in X^n$, with $i$-th component $v_i$. An \emph{$n
\times m$ matrix over $S$} is an $m$-length vector of $n$-length vectors over $S$ (its columns), with entry
$M_{ij}$ at row $i$ and column $j$; equivalently, an $n$-length vector of $m$-length vectors over $S$ (its
rows). The \emph{standard basis vector} $e_i \in S^n$ at index $i$ is the Kronecker $\delta$,
\begin{salign*}
   (e_i)_j =
      \begin{cases}
         1 & \text{if }i = j; \\
         0 & \text{otherwise}
      \end{cases}
\end{salign*}
\noindent The scalar \emph{dot product} of two same-length vectors is $u \cdot v = \sum_i u_i \cdot v_i$. Viewing vectors $w_1, \ldots, w_m \in S^n$ as the columns of an
$n \times m$ matrix $W$ (so $W_{ji} = (w_i)_j$), the iterated sum $\bigoplus_i w_i \in S^n$ is the row-sum of
$W$:
\begin{salign*}
   \left(\bigoplus_{i} w_i\right)_{\!j} = \sum_i W_{ji}.
\end{salign*}

\subsection{Basis decomposition}
\label{sec:matrix:basis-decomp}

Write $a \cdot w$ for the pointwise product of a scalar $a$ with a vector $w$, i.e.\ $(a \cdot w)_i = a
\cdot w_i$. Every $m$-length vector $v$ decomposes as an iterated sum of scaled basis vectors:
\begin{salign*}
   v = \bigoplus_{i} v_i \cdot e_i \tag{basis decomposition}
\end{salign*}
\noindent This follows from the symmetry $(e_i)_j = (e_j)_i$ of the Kronecker $\delta$ together with the
following lemma.

\begin{lemma}[Sum-of-basis]
\label{lem:matrix:sum-of-basis}
$\sum_j (e_i)_j \cdot v_j = v_i$.
\end{lemma}

\subsection{The matrix category}
\label{sec:matrix:category}

The product $N \comp M$ of an $n \times p$ matrix $N$ with a $p \times m$ matrix $M$ is the $n \times m$
matrix
\begin{salign*}
   (N \comp M)_{ij} = \sum_k N_{ik} \cdot M_{kj}.
\end{salign*}
\noindent This product is associative (using distributivity and the interchange law), with the $m \times m$
Kronecker $\delta$ matrix $I_{ij} = (e_i)_j$ as left and right unit (by \lemref{matrix:sum-of-basis}). The category $\Mat(S)$ accordingly has as objects
all natural numbers and as morphisms $m \to n$ all $n \times m$ matrices over $S$, with composition and
identities given by this product.

\subsubsection{Transpose}
\label{sec:matrix:transpose}

The \emph{transpose} $M^{\!\top} : n \to m$ of $M : m \to n$ has entries $(M^{\!\top})_{ji} = M_{ij}$. It is
involutive, $(M^{\!\top})^{\!\top} = M$, and interacts with composition as $(N \comp M)^{\!\top} = M^{\!\top}
\comp N^{\!\top}$ (both immediate from commutativity of $\cdot$), so $(-)^{\!\top}$ is a dagger on
$\Mat(S)$.

\subsection{$\CMon$-enrichment and biproducts}
\label{sec:matrix:cmon}

Pointwise addition $(M +_m N)_{ij} = M_{ij} + N_{ij}$ and the zero matrix $0_m$ equip each hom-set of $\Mat(S)$
with a commutative monoid; composition distributes pointwise over $+_m$, so that $\Mat(S)$ is
$\CMon$-enriched. Moreover, $0 : \Mat(S)$ is a zero object (both terminal and initial), and the sum $m + n$ of
natural numbers is the \emph{biproduct} of $m$ and $n$ in $\Mat(S)$, with injections and projections given by
the appropriate ``block'' matrices. $\Mat(S)$ is therefore a concrete semi-additive category.

\subsection{Additional structure on $S$}
\label{sec:matrix:additional-structure}

Further constructions become available when $S$ carries compatible order-theoretic structure.

If $+$ is the join of a bottom element $0$ — i.e.\ $(S, +, 0, \le)$ is a join-semilattice — then $\sum_i$
is iterated join (and $\bigoplus_i$ its pointwise lift), and basis decomposition extends to any
$\join$-preserving, scale-linear $f : S^m \to S^n$:
\begin{salign*}
   f(v) = \bigoplus_i v_i \cdot f(e_i) \tag{map basis decomposition}
\end{salign*}
\noindent so $f$ is determined by its values on the basis vectors $e_i$ (the columns of its matrix
representation).

If additionally $\cdot$ is the meet of a top element $1$ and distributes over $+$ — i.e.\ $(S, +, 0, \cdot,
1, \le)$ is a bounded distributive lattice — we have a \emph{disjointness} relation $a \disj b$ iff $a
\cdot b \le 0$, with pointwise lift to vectors.

If $S$ is furthermore a \emph{Boolean algebra}, equipped with a negation $\neg$ satisfying $1 \le a + \neg
a$ (excluded middle) and $a \cdot \neg a \le 0$ (non-contradiction), then $\neg$ is involutive and
disjointness reflects the order: $a \le b$ iff every $c$ disjoint from $b$ is disjoint from $a$.

\subsection{Example: $\Mat(\Two)$ for Boolean analysis}
\label{sec:matrix:two}

Take $S = \Two = (\{\bot, \top\}, \join, \bot, \meet, \top)$ with $+ := \join$, $\cdot := \meet$, $0 := \bot$, $1
:= \top$; this is a commutative (in fact Boolean) semiring. $\Mat(\Two)$ has as its morphisms $m \to n$ the $n
\times m$ Boolean matrices; matrix composition in the $(\Two, \join, \meet)$ semiring computes reachability
in the dependency graph encoded by a matrix (\secref{matrix:two:matrices}). This was the starting point for
the development and supplies the concrete semantics of interest for slicing-style analyses.

\subsection{Morphisms as dependency relations}
\label{sec:matrix:two:matrices}

At $S = \Two$ each entry $M_{ij}$ is a bit, and $M$ can be read as the \emph{dependency relation} between two
index sets: $M_{ij} = \top$ iff position $i$ in the codomain depends on position $j$ in the domain.
Composition in the $(\Two, \join, \meet)$ semiring is reachability:
\begin{salign*}
   (N \comp M)_{ik} = \bigjoin_j N_{ij} \meet M_{jk}
\end{salign*}
\noindent evaluates to $\top$ iff there is some intermediate position $j$ witnessing a chain $k \to j \to i$.
Basis decomposition specialises to
\begin{salign*}
   f(v) = \bigjoin_{j \in J(v)} f(e_j) \qquad (J(v) = \{j : v_j = \top\})
\end{salign*}
\noindent --- the ``live'' positions of $v$ select which columns of $M_f$ contribute.

\subsection{Relation to $\SemiLat$}
\label{sec:matrix:two:semilat}

$\Mat(\Two)$ is equivalent (as a $\CMon$-enriched category) to the full subcategory of $\SemiLat$ spanned by
the biproduct powers $\Two^n = \Two \biprod \cdots \biprod \Two$. The equivalence is witnessed by the
\emph{scalar iso} $\Two \cong \SemiLat(\Two, \Two)$, identifying $\bot$ with the zero morphism and $\top$ with
the identity; this lifts, via the semi-additive representation of morphisms as matrices, to the full
equivalence. This is captured in Agda by the module \texttt{matrix-embedding}, instantiated at $\Two$ in
\texttt{ho-model} (\secref{matrix:two:semilat}), which also exhibits that $\SemiLat(\Two, \Two)$ is commutative
under composition --- a fact which transports across the iso from commutativity of $(\Two, \meet)$, removing
the need for a direct case analysis.

\subsection{Two embeddings (to be generalised)}
\label{sec:matrix:two:embeddings}

The following material is the Boolean-specific account of the two embeddings, to be folded into a generic
account (valid in the Boolean algebra setting of \secref{matrix:additional-structure}) in a later pass.

We now treat $f: \Two^m \to \Two^n$ as the \emph{backward} direction of an analysis --- a map that propagates
output demand back to input demand, and consider two contravariant embeddings: into $\BoolConj$ (a subcategory
of $\LatGal$), and into $\LatGal$ directly. The two embeddings pair $f$ with a forward-direction counterpart
in two different ways:

\begin{enumerate}
\item \emph{Conjugate pair.} The transpose $f^{\!\top}$ pairs with $f$ to form a conjugate pair $(f^{\!\top},
f)$ in $\BoolConj$ (\secref{categories-with-biproducts:boolconj}): both maps are $\join$-preserving, linked by
the disjointness condition
\begin{salign*}
   y \disj f^{\!\top}(x) \iff f(y) \disj x
\end{salign*}

\item \emph{Galois connection.} The right adjoint $f_*: \Two^n \to \Two^m$ (meet-preserving) pairs with $f$ to
form a Galois connection $f \dashv f_*$ in $\LatGal$ (\secref{categories-with-biproducts:latgal}), with
adjointness condition
\begin{salign*}
   y \le f_*(x) \iff f(y) \le x
\end{salign*}
\noindent Concretely, $f_*(x)_i = \bigmeet_j (M_f[i,j] \to x_j)$, where $a \to b = \neg a \join b$ is
Boolean implication.
\end{enumerate}

\subsection{De Morgan duality}
\label{sec:matrix:two:de-morgan}

The two backward maps are De Morgan duals:
\begin{salign*}
   \neg\,f^{\!\top}(x) = f_*(\neg x)
\end{salign*}
\noindent This is the Boolean-algebra identity relating the transpose (in the $(\join, \meet)$-semiring) to
the adjoint (in the dual $(\meet, \join)$-semiring). Combined with the Boolean-algebra fact that $a \disj b
\iff a \le \neg b$, we obtain the following chain:
\begin{salign*}
   y \disj f^{\!\top}(x)
   & \iff y \le \neg\,f^{\!\top}(x) \\
   & \iff y \le f_*(\neg x) \\
   & \iff f(y) \le \neg x \\
   & \iff f(y) \disj x
\end{salign*}
\noindent This (and an analogous argument for the Galois connection) exhibits the conjugate and Galois
embeddings as mutually derivable given the Boolean structure on $\Two^n$. This is specific to the Boolean
case. In a general Heyting algebra the equivalence $a \disj b \iff a \le \neg b$ still holds (taking $\neg$ as
the pseudo-complement), but the De Morgan identity $\neg f^{\!\top} = f_* \comp \neg$ relies on
double-negation $\neg \comp \neg = \id$, which is Boolean-specific.

\subsection{Connection to forward- and reverse-mode AD}
\label{sec:matrix:two:ad}

\begin{center}
\small
\begin{tabular}{ll}
 \textbf{Real-number AD in $\FinVect_\RR$} & \textbf{Boolean AD in $\Mat(\Two)$} \\
 Linear map $\pushf{\phi}_x: \RR^m \linearto \RR^n$ (forward mode) & $\join$-preserving transpose $f^{\!\top}: \Two^m \to \Two^n$ (demanded by) \\
 Linear transpose $\pullf{\phi}_x: \RR^n \linearto \RR^m$ (reverse mode) & Given $\join$-preserving $f: \Two^n
 \to \Two^m$ (demands) \\
 Matrix multiplication in $(\RR, +, \cdot)$ semiring & Matrix multiplication in $(\Two, \join,
 \meet)$ semiring
\end{tabular}
\end{center}
